\cgeentry{cge_results}{Report simulation results}

\cgesec{Description}

\cmd{cge_results} displays a summary of a simulation: real value added, the
exchange rate, the CPI, investment, imports, exports, government savings,
factor prices and activity output, as percentage changes from the base; the
equivalent variation of each household; and the checks of Walras's law and of
the GDP identity. With \opt{}{clear} or \opt{sav}{ing()}, it writes every
variable of the solution as a dataset, with one observation per variable
element.

\cgesec{Quick start}

Display the summary of simulation \cmd{notariff}
\begin{stcode}
cge_results, sim(notariff)
\end{stcode}
Load every variable of the simulation as data, without the summary
\begin{stcode}
cge_results, sim(notariff) clear nodisplay
\end{stcode}
Save the results of the last simulation as a Stata dataset
\begin{stcode}
cge_results, saving(results, replace)
\end{stcode}

\cgesec{Syntax}

\begin{stsyntax}
\cmdb{cge_results} [\texttt{,} \argu{options}]
\end{stsyntax}

\begin{opttable}
\optgroup{Main}
\optarg{}{sim}{name} & simulation; default the last one stored\\
\optgroup{Data}
\opt{}{clear} & replace the data in memory with the results\\
\optarg{sav}{ing}{filename\textup{[\texttt{, replace}]}} & save the results as a Stata dataset\\
\optgroup{Reporting}
\opt{nod}{isplay} & do not display the summary\\
\end{opttable}

\cgesec{Options}

\begin{optdesc}
\item[\optarg{}{sim}{name}] selects the simulation. The default is the last
simulation stored, which is the one most recently added: solving again under
an existing name replaces that simulation where it stands.

\item[\opt{}{clear}] replaces the data in memory with the results of the
simulation. The data in memory need not have been saved.

\item[\optarg{sav}{ing}{filename\textup{[\texttt{, replace}]}}] saves
the results as a Stata dataset, leaving the data in memory unchanged;
\cmd{replace} allows an existing file to be overwritten.

\item[\opt{nod}{isplay}] suppresses the summary. The summary also stores the
results in \cmd{r()}, so nothing is stored with \opt{nod}{isplay}.
\end{optdesc}

\cgesec{Remarks and examples}

\cmd{cge_results} is the quick look at one simulation, and the way to every
variable of its solution. \manref{cge_report} computes the standard report
measures, such as real GDP and revenue by tax, for several simulations side by
side.

\cgesubsec{Example 1: the summary}

In the synthetic SAM, mapped and calibrated as in \manref{cge_solve}, every
import tariff is removed and the model solved as simulation \cmd{notariff}:
\stlog{res_show}
The heading gives the closure under which the simulation was solved. The
percentage changes are relative to the absolute value of the base, so a fall is
negative even when the base is negative: here government savings, 3 in the
base, fall by 271 per cent to a deficit of about 5. The household gains, with
an equivalent variation of 6.1, 2.7 per cent of its base consumption. The two
checks are zero to the precision of the solution: the slack of the
savings--investment balance, \cmd{WALRAS}, and the gap between nominal GDP
measured from expenditure and from income.

The equivalent variations are returned in a row vector, one column per
household:
\stlog{res_return}

\cgesubsec{Example 2: the results as data}

With \opt{}{clear}, every variable element of the solution becomes an
observation, with its base and simulated values:
\stlog{res_data}
The variables are \cmd{variable}, the name of the model variable as in
\secref{sec:variables}; \cmd{element1} and \cmd{element2}, its set elements,
empty for a variable without index; \cmd{base}, the base value; and
\cmd{value}, the value in the simulation. Every variable is included, fixed or
free, with the rents of constraints, the slack \cmd{WALRAS} and the user
variables of \manref{cge_variable}; elements that do not exist in the model,
such as imports of a commodity that is not imported, are left out. Percentage
changes are then a line away:
\stlog{res_pct}
The 15 missing values are the elements whose base value is zero.

\cgesubsec{Example 3: saving the results}

\opt{sav}{ing()} writes the same dataset to a file and leaves the data in
memory unchanged:
\stlog{res_saving}

\cgesec{Stored results}

\cmd{cge_results} stores the following in \cmd{r()}, unless
\opt{nod}{isplay} is specified:
\begin{storedres}
\resgroup{Macros}
\res{r(sim)} & name of the simulation\\
\resgroup{Matrices}
\res{r(ev)} & equivalent variation of each household, $1 \times H$\\
\end{storedres}

\cgesec{Methods and formulas}

Let $x^0$ be a variable at the base and $x$ in the simulation. The summary
shows $100\,(x - x^0)/|x^0|$, missing when $x^0 = 0$. Its rows are real value
added, $\sum_a \V{PVA}^0_a\,\V{QVA}_a$ less the commodities in nests of value
added at base prices (\secref{sec:nests}); $\V{EXR}$; $\V{CPI}$; $\V{IADJ}$;
imports, $\sum_c \V{pwm}_c\,\V{QM}_c$, and exports,
$\sum_c \V{pwe}_c\,(\V{QE}_c + \V{QER}_c)$, both at the world prices of the
simulation for the base and for the solution alike; $\V{GSAV}$; $\V{SHIADJ}$,
when the model has a health insurance fund; $\V{W}_f$; and $\V{QA}_a$.

The equivalent variation of household $h$ is the change in spending at base
prices that gives its new utility,
\begin{equation*}
\V{EV}_h = \sum_c P^0_c\,\gamma_{ch}
  + \Bigl(\V{EH}_h - \sum_c P_c\,\gamma_{ch}\Bigr)
    \prod_c \Bigl(\frac{P^0_c}{P_c}\Bigr)^{\beta_{ch}} - \V{EH}^0_h,
\end{equation*}
where the sums and the product run over marketed commodities, at the prices
that final users pay, $\V{PH}_c$, and over home-consumed goods, at the activity
prices $\V{PA}_a$, with the shares and subsistence quantities of home
consumption; $\gamma_{ch} = 0$ with Cobb--Douglas demand. It is shown in per
cent of base consumption, $\V{EH}^0_h$. The GDP identity gap is nominal GDP
measured from expenditure less nominal GDP measured from income, zero at any
solution. \Chapref{chap:reporting} gives the measures, the compensating
variation and the two sides of GDP (\secref{sec:welfare}).

\cgesec{Also see}

\begin{alsosee}
\seeentry{cge_report}{Standard report tables for one or more simulations}
\seeentry{cge_solve}{Solve the model}
\seeentry{cge_check}{Run model validation checks}
\seeentry{cge_parameters}{The calibrated parameters as a dataset}
\seechapter{chap:reporting}{Reporting measures and welfare}
\end{alsosee}
